\documentclass[12pt]{article}
\usepackage{anyfontsize}
\usepackage[a4paper, margin=2cm]{geometry}
\usepackage{polski}
\usepackage{tabto}
\usepackage{enumitem}
\usepackage{amsmath}
\usepackage{amssymb}
\usepackage{multirow}
\usepackage{multicol}
\usepackage{setspace}
\usepackage{pdfpages}



\usepackage{tabularx}
\newcolumntype{C}{>{\centering\arraybackslash}X}
\newcolumntype{L}{>{\raggedleft\arraybackslash}X}
\newcolumntype{R}{>{\raggedright\arraybackslash}X}
\newcommand{\centerY}[2]{\multirow{#1}{*}{#2}}

\usepackage{wrapfig}

\usepackage{hyperref}
\hypersetup{
    colorlinks = true,
    urlcolor=blue,
    linkcolor= black
}

\usepackage{chngcntr}
\counterwithin{figure}{section}
\counterwithin{table}{section}
\numberwithin{equation}{section}

\usepackage{graphicx}
\graphicspath{{./Img/}}

\usepackage{csvsimple}
\usepackage{pgfplots}
\usepackage{pgfplotstable}
\pgfplotsset{compat= newest}


\usepackage{titlesec}
\titlelabel{\thetitle.\quad}
% \AddToHook{cmd/section/before}{\clearpage}

\usepackage[european, american currents, americanvoltages, RPvoltages, cute inductor]{circuitikz}
\usepackage{tikz}
\usetikzlibrary{shapes.geometric}
\ctikzset{
    logic ports=ieee,
    logic ports/scale=0.7,
}


\usepackage{listings}
\lstset{
literate=%
    {ą}{{\k{a}}}1
    {Ą}{{\k{A}}}1
    {ć}{{\'c}}1
    {Ć}{{\'{C}}}1
    {ę}{{\k{e}}}1
    {Ę}{{\k{E}}}1
    {ł}{{\l{}}}1
    {Ł}{{\L{}}}1
    {ń}{{\'n}}1
    {Ń}{{\'N}}1
    {ó}{{\'o}}1
    {Ó}{{\'O}}1
    {ś}{{\'s}}1
    {Ś}{{\'S}}1
    {ż}{{\.z}}1
    {Ż}{{\.Z}}1
    {ź}{{\'z}}1
    {Ź}{{\'Z}}1
}


\title{
    \includegraphics[width = 0.3\textwidth]{agh_logo.jpg}\\
    \textbf{Akademia Górniczo-Hutnicza w Krakowie}\\
    Wydział Informatyki, Elektroniki i  Telekomunikacji\\\vspace{2cm}
    \textbf{Raport z projektu}\\
    Analizator stanów logicznych
}
\author{
    \begin{tabularx}{\textwidth}{l l}
    Autor: &Łukasz Przystupa, Krzysztof Płonka, Paweł Olbrych\\
    Kierunek studiów: & Elektronika i Telekomunikacja\\
    Przedmiot:& Sensory w Aplikacjach Wbudowanych
    \end{tabularx}
}
\date{\vspace{2cm}\today}

\usepackage{titling}
\renewcommand\maketitlehooka{\null\mbox{}\vfill}
\renewcommand\maketitlehookd{\vfill\null}

\begin{document}
    \begin{titlepage}
        \maketitle
        \thispagestyle{empty}
    \end{titlepage}

    \pagenumbering{Roman}
        \include{shorts.tex}
        \tableofcontents
        \newpage
    \pagenumbering{arabic}

    \section{Opis i założenia projektu}
    Celem projektu jest zaprojektowanie oraz zbudowanie systemu analizatora stanów logicznych.
    System analizatora składa się z 2 części:
    \begin{itemize}
        \item Układu analizatora - dedykowane PCB zawierające min. mikrokontroler
        Raspberry Pi Pico W, przetwornik ADC, stabilne źródło napięcia odniesienia(bandgap),
        układ filtracji sygnału analogowego(LPF) oraz interfejs wejściowy wraz z zabezpieczeniami
        chroniącymi układ przed przypadkowymi pomyłkami użytkownika.
        \item Aplikacji komputerowej - interfejs użytkownika odpowiadający za intuicyjną komunikację 
        między użytkownikiem a urządzeniem. Aplikacja komputerowa stworzona zostanie z wykorzystaniem
        framework-u GTK4 oraz języka C. GUI dodatkowo będzie wykorzystywać mechanizmy wielowątkowości w celu
        zapewnienia odpowiedniej kultury pracy z aplikacją. 
    \end{itemize}

\subsection{Założenia sprzętowe}
    \begin{enumerate}
        \item Analizator stanów logicznych, jako układ pomiarowy nie powinien obciążać urządzenia badanego dlatego
        system pomiarowy analizatora będzie mieć możliwość ustawienia wyprowadzeń w stan wysokiej impedancji.
        \item Ukłąd powinien charakteryzować się co najmniej dobrymi parametrami takimi jak:
        szybkość pracy, rozdzielczość sygnałów, dokładność oraz precyzja pomiaru.
        \item System powinien umożliwiać pomiary sygnałów których częstotliwość bitowa będzie wynosić co najmniej
        100kB/s %tak se podałem te 100 bo nie wiem ile ma mieć 
    \end{enumerate}

    section{Założenia oprogramowania}
    \begin{enumerate}
        \item Aplikacja komputerowa powinna stanowić wygodny oraz intuicyjny interfejs między użytkownikiem
        a urządzeniem.
        \item Oprogramowanie mikrokontroler w maksymalnym stopniu powinno opierać się na zasobach
        hardware-owych takich jak DMA, Timery, akceleratory IO(PIO) itp. .
    \end{enumerate}
    
% TO bym dał potem a w wstępie dał tylko założenia
% \subsection{Założenia oprogramowania}
%     \begin{enumerate}
%         \item Projekt zakłada stworzenie dwóch aplikacji.
%         Pierwszą z nich będzie oprogramowanie mikrokontrolera, będącego sercem części hardwarowej.
%         Drugą zaś aplikacja graficzna pozwalająca wyświetlić użytkownikowi dane w wygodny sposób.
%     \end{enumerate}

% \subsubsection{Oprogramowanie mikroprocesora}
%     \begin{enumerate}
%         \item Mikrokontroler w maksymalnym stopniu powinien opierać się na zasobach hardwareowych, takich jak DMA czy akceleratory IO.
%         Natomiast komunikacja z komputerem, powinna być możliwie prosta - przykładowo wykorzystując protokół UART.
%         A w znacznie lepszym wykorzystaniem byłoby standard CDC (Communication Device Class) -- wirtualne połączenie przez port szeregowy,
%         dzięki czemu można przesyłać dane z maksymalną przepustowością \textit{USB 2.0}.
%     \end{enumerate}


\subsection{Podział obowiązków}
    \begin{itemize}
        \item Łukasz Przystupa  - napisanie aplikacji graficznej oraz oprogramowanie Raspberry PI PICO.
        \item Krzysztof Płonka  - stworzenie interfejscu do komunikacji bezprzewodowej(Wi-Fi),
        biblioteki do układu ADC oraz (opcjonalne) ... . % TODO
        \item Paweł Olbrych     - zaprojektowanie PCB.
    \end{itemize}
    \section{Opis układów elektronicznych wykorzystanych w projekcie}

Do najważniejszych elementów elektronicznych wykorzystanych 
podczas projektowania analizatora należą min.:
 
\begin{enumerate}
    \item Raspberry Pi Pico W - platforma firmy Raspberry PI.
    % układ RP2040 oraz układ CYW43439. 
        Płytka wyposażona jest w mikrokontroler z rodziny ARM Cortex M0+, zaprojektowany przez firmę Raspberry Pi Foundation.
        Do komunikacji bezprzewodowej mikrokontroler wykorzystuje zewnętrzny chip firmy Inineon o nazwie CYW43439.
        Układ ten pozwala na komunikację z wykorzystanie zarówno BLE jak i WIFI.\\
        Jednym z peryferiów mikroprocesora RP2040, są dwa układy PIO.
        Moduły te pozwalają na moduły te pozwalają na bardzo szybkie, niezależne od głównego programu zarządzanie układami IO.
        A w połączeniu z układami DMA, możliwe jest próbkowanie i zapisywanie sygnałów wejściowych z częstotliwością pracy głównego rdzenia ($\approx 200MHz$).
    \item ADC ADS1115 - 
        układ przetwornika analogowo-cyfrowego typu Sigma-Delta o rozdzielczości 16 bitów.
        Układ może działać z maksymalną prędkością 860 sampli,
        a dzięki wykorzystaniu wbudowanego multipleksera może przełączać się między 4 wejściami single ended lub 2 różnicowymi.
        Przetwornik został wyposażony w prosty interfejs komunikacyjny I2C za pomocą, którego możliwa jest jego konfiguracja oraz przesył danych.
    \item Stabilizator liniowy LDO ... TODO
    \item Zestaw 4 filtrów analogowych ... TODO  
    \item Diody zabezpieczające - diody Zenera na napięcie 3.6V zabezpieczające stopień wejściowy analizatora przed przypadkowymi przepięciami i zwarciami.
    \item ... TODO
\end{enumerate}


    \section{Opis schematu}
        TODO


    \vfill
    \begin{flushright}
        Łukasz Przystupa,\\
        Krzysztof Płonka,\\
        Paweł Olbrych
    \end{flushright}
\end{document}
